\section{引言}

随着在轨服务、空间碎片清理和失效航天器维护需求的增长，空间机械臂已成为执行非合作目标接近、捕获与操作任务的重要装备\cite{liu2021review}。与地面固定基座机械臂不同，捕获前空间机械臂通常处于自由漂浮状态，基座不直接受控，机械臂关节运动会通过系统动量守恒关系诱导基座线速度和角速度变化，进而改变末端执行器的实际运动。严宇新针对在轨捕获任务指出，空间机械臂动力学模型具有高度非线性和强耦合特性，并且轨迹跟踪过程中还会受到模型参数不确定、未知时变扰动和执行器故障等因素影响\cite{yan}。因此，如何在基座不直接施加控制力矩的条件下实现高精度、可预测收敛时间的轨迹跟踪，是自由漂浮空间机械臂控制中的关键问题。

自由漂浮空间机械臂的控制难点首先来自建模层面的非完整约束。在线动量和角动量初值为零的自由漂浮模式下，基座速度可由机械臂关节速度唯一诱导，末端速度不再由固定基座雅可比矩阵决定，而应由包含基座反作用运动的广义雅可比矩阵描述。

Umetani 和 Yoshida 提出的广义雅可比矩阵为自由漂浮空间机械臂的速度级建模和分解运动速度控制奠定了经典基础\cite{umetani1989}；Vafa 和 Dubowsky 的虚拟机械臂方法从动量守恒角度系统刻画了空间机械臂运动学与动力学耦合\cite{vafa1990}；Papadopoulos 和 Dubowsky 进一步揭示了自由漂浮空间机械臂中由动力学耦合引起的动态奇异问题，并系统总结了自由漂浮与自由飞行两类空间机器人的运动学与动力学特性\cite{papadopoulos1993,dubowsky1993}。在动力学建模与控制层面，Parlaktuna 和 Ozkan 提出了基于动力学等效机械臂的自适应计算力矩控制方案\cite{parlaktuna2004}；Xu 等进一步设计了自适应零反作用运动控制器，在跟踪末端轨迹的同时主动抑制基座姿态扰动\cite{xu2016}。严宇新论文中由动量守恒方程推导出基座--机械臂雅可比矩阵 \(J_{bm}\)，并进一步得到 \(J_g=J_bJ_{bm}+J_m\) 的自由漂浮运动学关系；在动力学层面，通过消去基座加速度，可将系统写成 \(H_g(q)\ddot q+C_g(q,\dot q)\dot q=\tau_m\) 的等效关节空间形式\cite{yan}。这一约化形式为把自由漂浮系统转化为关节空间控制问题提供了基础。

已有研究围绕轨迹规划、约束处理和鲁棒跟踪展开。针对自由漂浮模式下的运动学冗余、动力学奇异和关节速度约束，Wang 等采用粒子群优化生成满足约束的关节轨迹\cite{wang2015}；宁昕和武耀发针对传统控制方法难以同时处理自由漂浮空间机器人轨迹跟踪约束的问题，基于拉格朗日动力学模型构造扩展状态空间模型，并利用模型预测控制实现对末端期望位置和速度的跟踪\cite{ning2019}；针对捕获前反作用抑制和参数不确定问题，也有研究将反作用零空间规划、时延估计和自适应滑模控制结合，以降低对精确动力学模型的依赖\cite{yu2022}。Yao 基于障碍李雅普诺夫函数和径向基神经网络，为自由飞行空间机械臂设计了全状态和输出反馈自适应预设性能控制器，确保位置跟踪误差在参数不确定、外部扰动和输入饱和条件下始终保持在预定性能边界内\cite{yao2021}。这些方法在跟踪精度和鲁棒性方面各有侧重，但关于跟踪误差的收敛时间通常依赖于控制器参数与初始状态的复杂关系，难以在设计阶段独立指定。

近年来，固定时间（fixed-time）、预定义时间（predefined-time）和指定时间（prescribed-time）控制进一步受到关注，其共同优势在于可在设计阶段给出与初始误差弱相关或无关的收敛时间估计。在非线性系统的固定时间控制方面，Polyakov 提出了基于非连续齐次性的固定时间镇定框架\cite{polyakov2012}；Zuo 等进一步将固定时间方法推广至多智能体系统一致性控制\cite{zuo2019}。对于机器人系统，Su 等利用非奇异终端滑模实现了机械臂的固定时间轨迹跟踪控制\cite{su2020}；相关研究已将预定义时间稳定思想用于自由飞行空间机器人 SE(3) 位姿跟踪\cite{xuPredefinedTimeTrackingControl2024}，也将固定时间容错控制用于含执行器故障和性能约束的自由漂浮空间机械臂\cite{li2025}。然而，上述固定时间或预定义时间方法大多依赖非光滑终端滑模面或分数幂反馈项，在收敛时刻附近存在控制输入跳变或增益奇异问题，对工程实现构成挑战。

离散周期延迟反馈（delayed feedback control, DFC）概念最早由 Pyragas 提出，用于镇定混沌系统中的不稳定周期轨道\cite{pyragas1992}。其基本思想是通过系统当前状态与延迟一个周期长度的状态之差构造反馈信号，在不改变目标轨道的前提下实现镇定。延拓至连续时间框架后，Zhou、Michiels 和 Chen 提出的光滑周期延迟反馈（smooth periodic delayed feedback, PDF）方法表明，对于可控线性时滞系统，若预设收敛时间满足相应可解条件，可通过周期延迟反馈实现固定时间镇定；其反馈增益可选为连续、连续可微甚至光滑函数，从而避免传统指定时间高增益方法在终端时刻出现的增益奇异\cite{zhouFixedtimeStabilizationLinear2022}。该方法的核心不是构造非光滑终端滑模面，而是利用周期延迟系统的有限维单周期状态转移矩阵（又称单值矩阵），并通过使该矩阵零化或幂零来获得固定时间收敛。该方法在线性系统和线性化模型框架下具有理论简洁性和优越的时域性能，但其适用性需要被控系统首先转换为线性可控形式。

需要说明的是，固定时间收敛结论从非线性原系统到线性误差通道的迁移依赖于扰动补偿的精度。预设时间扰动观测器（prescribed-time disturbance observer, PTDO）是达成这一目的的重要工具。已有研究表明，通过时变增益构造可保证观测误差在预设时间内精确收敛至零\cite{jiangPrescribedtimeDisturbanceObserver2024}。在机器人控制领域，扰动观测器结合计算力矩补偿的方法已得到广泛验证，但将其与固定时间反馈结构显式衔接并给出联合预设时间收敛保证的相关工作尚不充分。

基于以上背景，本文采用"扰动观测--动力学补偿--误差线性化--周期延迟反馈"的结构：先依据动量守恒关系将捕获前自由漂浮空间机械臂约化为等效关节空间模型，再通过预设时间扰动观测器得到集总扰动估计 \(\hat d(t)\)，随后利用计算力矩补偿得到从辅助加速度到跟踪误差的线性可控通道，最后在该误差通道中嵌入光滑周期延迟反馈。这样既保留 PDF 方法的固定时间镇定机理，又使其作用对象与原理论假设保持一致。

本文主要工作如下：第一，基于自由漂浮空间机械臂的动量守恒约束，给出从基座反作用运动到广义雅可比矩阵、等效关节空间动力学的推导链条；第二，引入预设时间扰动观测补偿机制，说明 \(\hat d(t)\) 可由等效加速度扰动观测得到，并给出观测--PDF 两阶段设计时序；第三，构造面向 PDF 设计的误差通道规范化方法，将关节轨迹跟踪问题转化为线性可控误差系统镇定问题；第四，结合 PDF 理论中的 \(R_h(t)\)、可控 Gramian 和单周期状态转移矩阵零化条件，给出关节力矩控制律及其固定时间收敛条件；第五，基于七自由度空间机械臂参数搭建 MATLAB 仿真，验证所提控制结构的可运行性和有界扰动下的跟踪性能。
